\chapter{Os órgãos em atividade}\label{ch:g7:organs-at-work}

Suba correndo três lances de escada e faça o balanço: \omterm{def:g1:my-body:parts}{pernas} ardendo,
\omterm{def:g4:heart-and-blood:heart}{coração} martelando, \omterm{def:g7:what-is-respiration:respiration}{respiração} disparada --- e, curiosamente, a \omterm{def:g1:the-five-senses:sense}{pele}
avermelhada e quente. Ninguém mandou essas mudanças acontecerem; elas
se mandaram sozinhas, e em proporção perfeita. Este capítulo mede o
que um órgão em atividade de fato consome, e como o corpo reorganiza
as entregas dele no instante em que o trabalho começa.

\section{O que um músculo em atividade consome}

\begin{proposition}[A conta do músculo]\label{prop:g7:organs-at-work:bill}
Um \omterm{def:g3:bones-and-muscles:muscle}{músculo} em atividade gasta, e o gasto dele pode ser medido.
Comparado com o mesmo \omterm{def:g3:bones-and-muscles:muscle}{músculo} em repouso, um \omterm{def:g3:bones-and-muscles:muscle}{músculo} trabalhando duro:
\begin{itemize}
  \item retira muito mais \emph{\omterm{def:g7:what-is-respiration:respiration}{oxigênio}} do \omterm{prop:g4:journey-of-food:absorb}{sangue} que passa por ele
        --- várias vezes a retirada de repouso;
  \item retira muito mais \emph{glicose}\index{glicose} --- o
        \omterm{def:g4:journey-of-food:digestion}{nutriente} açucarado pronto do \omterm{prop:g4:journey-of-food:absorb}{sangue}, o combustível do dia a
        dia dele;
  \item libera na mesma proporção mais \emph{\omterm{def:g7:what-is-respiration:respiration}{gás carbônico}} --- e
        \emph{calor}.
\end{itemize}
Isso é a queima celular da \cref{prop:g7:what-is-respiration:power},
lida nos medidores de um órgão: combustível e \omterm{def:g7:what-is-respiration:respiration}{oxigênio} para dentro,
trabalho, gás residual e calor para fora.
\end{proposition}
\admitted

\begin{example}[Medindo através de um músculo]\label{ex:g7:organs-at-work:measure}
Os fisiologistas comparam o \omterm{prop:g4:journey-of-food:absorb}{sangue} que \emph{entra} num \omterm{def:g3:bones-and-muscles:muscle}{músculo} com o
\omterm{prop:g4:journey-of-food:absorb}{sangue} que \emph{sai} dele. Em repouso, o \omterm{prop:g4:journey-of-food:absorb}{sangue} que sai entregou uma
parcela modesta do \omterm{def:g7:what-is-respiration:respiration}{oxigênio} e da \omterm{prop:g7:organs-at-work:bill}{glicose} dele. No trabalho pesado, a
mesma comparação mostra o \omterm{prop:g4:journey-of-food:absorb}{sangue} que sai raspado --- com quase todo o
\omterm{def:g7:what-is-respiration:respiration}{oxigênio} entregável retirado --- enquanto o \omterm{def:g7:what-is-respiration:respiration}{gás carbônico} deu um salto.
O mesmo órgão, os mesmos medidores: só o trabalho mudou, e a conta o
acompanhou.
\end{example}

\begin{remark}[O calor é o recibo]\label{rem:g7:organs-at-work:warmth}
O calor do corpo em atividade não é acidente: queimar combustível
sempre produz calor, e o \omterm{def:g3:bones-and-muscles:muscle}{músculo} é a sala das caldeiras do corpo.
Tremer de frio é esse truque acionado de propósito --- \omterm{def:g3:bones-and-muscles:muscle}{músculos} postos
a trabalhar para produzir só calor (a despesa fixa do
\cref{ex:g3:balanced-diet:rest}, dobrada por um instante). A \omterm{def:g1:the-five-senses:sense}{pele}
avermelhada da corrida é o sistema de resfriamento levando embora o
calor da caldeira --- assunto da próxima seção.
\end{remark}

\section{Entregas sob demanda}

\begin{proposition}[O corpo reorganiza o abastecimento dele]\label{prop:g7:organs-at-work:supply}
No instante em que o trabalho começa, as linhas de abastecimento se
reorganizam para pagar a conta (o
\cref{ex:g4:heart-and-blood:effort} anunciou isso; aqui está o
despacho completo):
\begin{enumerate}
  \item os \emph{movimentos respiratórios} ficam mais fundos e mais
        rápidos --- recarregando o \omterm{prop:g4:journey-of-food:absorb}{sangue} mais depressa nos \omterm{def:g7:how-animals-breathe:four}{pulmões}
        (\cref{ex:g4:breathing:running});
  \item o \emph{\omterm{def:g4:heart-and-blood:heart}{coração}} bate mais rápido e com mais força --- mais
        rodadas por minuto
        (\cref{prop:g4:heart-and-blood:pulse});
  \item o \omterm{prop:g4:journey-of-food:absorb}{sangue} é \emph{redirecionado}: os \omterm{def:g4:heart-and-blood:vessels}{vasos} que servem os
        \omterm{def:g3:bones-and-muscles:muscle}{músculos} em atividade (e a \omterm{def:g1:the-five-senses:sense}{pele} que resfria) se alargam,
        enquanto os do \omterm{def:g4:journey-of-food:tract}{tubo digestivo} em repouso se estreitam --- o
        mesmo \omterm{prop:g4:journey-of-food:absorb}{sangue}, redirecionado para o andar que está gastando.
\end{enumerate}
Os três acontecem automaticamente conforme o trabalho, conduzidos
pelos serviços que o \omterm{def:g5:brain-in-command:system}{cérebro} presta sem que ninguém peça
(\cref{rem:g5:brain-in-command:more}).
\end{proposition}
\admitted

\begin{example}[O despacho, sentido por dentro]\label{ex:g7:organs-at-work:felt}
O balanço da subida de escada, explicado linha por linha: \omterm{def:g7:what-is-respiration:respiration}{respiração}
disparada --- recarregando; \omterm{def:g4:heart-and-blood:heart}{coração} martelando --- rodadas mais
rápidas; \omterm{def:g1:my-body:parts}{pernas} ardendo --- a conta dos \omterm{def:g3:bones-and-muscles:muscle}{músculos} sendo paga no limite;
\omterm{def:g1:the-five-senses:sense}{pele} quente e avermelhada --- o calor da caldeira levado até a
superfície de resfriamento. Até a clássica pontada e a moleza depois de
uma refeição grande se encaixam: a \omterm{def:g4:journey-of-food:digestion}{digestão} e a corrida disputam o
\omterm{prop:g4:journey-of-food:absorb}{sangue} redirecionado, e é por isso que o
\cref{met:g4:journey-of-food:habits} desaconselhava correr de barriga
cheia.
\end{example}

\begin{omfigure}
\begin{tikzpicture}
\begin{axis}[omaxis,
    width=9.6cm, height=5.8cm,
    xlabel={minutos}, ylabel={batidas por minuto},
    xmin=0, xmax=13.6, ymin=40, ymax=200,
    xtick={0,2,4,6,8,10,12},
    ytick={60,100,140,180},
    ]
  \addplot[very thick, omThm, smooth] coordinates {
    (0,72) (1,74) (2,120) (3,150) (4,158) (5,160)
    (6,120) (7,100) (8,88) (9,80) (10,76) (11,74) (12,73)
  };
  \draw[dashed, omExo] (axis cs:2,40) -- (axis cs:2,180);
  \draw[dashed, omExo] (axis cs:5,40) -- (axis cs:5,180);
  \node[font=\footnotesize] at (axis cs:1.0,52) {repouso};
  \node[font=\footnotesize] at (axis cs:3.5,52) {esforço};
  \node[font=\footnotesize] at (axis cs:8.6,52) {recuperação};
\end{axis}
\end{tikzpicture}

{\small Um \omterm{prop:g4:heart-and-blood:pulse}{pulso} ao longo de três minutos de esforço pesado: a subida
junto com o trabalho, o platô perto do máximo e a volta lenta ---
\omterm{def:g7:organs-at-work:recovery}{recuperação} --- enquanto o corpo acerta as contas dele.}
\end{omfigure}

\begin{definition}[Recuperação]\label{def:g7:organs-at-work:recovery}
Quando o esforço para, o \omterm{prop:g4:heart-and-blood:pulse}{pulso} e a \omterm{def:g7:what-is-respiration:respiration}{respiração} não param: eles caem só
aos poucos --- é a fase de \emph{recuperação}\index{recuperação}. O
corpo está acertando dívidas: recarregando as reservas de \omterm{def:g7:what-is-respiration:respiration}{oxigênio},
eliminando os resíduos do esforço, soltando o calor acumulado. Quanto
mais em forma está o corpo, mais rápida é a recuperação dele --- e é
por isso que o tempo de recuperação é, ele mesmo, uma medida de
condicionamento.
\end{definition}

\section{O treino e os limites da máquina}

\begin{proposition}[O treino levanta o teto]\label{prop:g7:organs-at-work:training}
Exercitada com regularidade, a cadeia de abastecimento inteira fica
mais forte (\cref{rem:g4:heart-and-blood:rest},
\cref{met:g4:breathing:care}): o \omterm{def:g4:heart-and-blood:heart}{coração} move mais \omterm{prop:g4:journey-of-food:absorb}{sangue} por batida,
os \omterm{def:g3:bones-and-muscles:muscle}{músculos} respiratórios ficam mais fundos, os próprios \omterm{def:g3:bones-and-muscles:muscle}{músculos}
crescem e guardam mais combustível, e os \omterm{def:g4:heart-and-blood:vessels}{vasos} finos deles se
multiplicam. Resultado: o mesmo esforço custa a um corpo treinado um
\omterm{prop:g4:heart-and-blood:pulse}{pulso} mais baixo, e o teto dele --- o esforço mais pesado que ele
consegue pagar --- sobe. A máquina é reconstruída pelo próprio uso.
\end{proposition}
\admitted

\begin{example}[Dois corredores, uma ladeira]\label{ex:g7:organs-at-work:runners}
Um sem treino e um treinado sobem a mesma ladeira. Sem treino: \omterm{prop:g4:heart-and-blood:pulse}{pulso} a
$180$ na metade, ofegante, obrigado a andar --- as linhas de
abastecimento não conseguem pagar a conta, e a dívida manda parar.
Treinado: \omterm{prop:g4:heart-and-blood:pulse}{pulso} a $150$ no topo, conseguindo falar, \omterm{def:g7:organs-at-work:recovery}{recuperação} em
três minutos. A mesma ladeira, a mesma conta --- tetos diferentes,
construídos por meses de ladeiras exatamente assim, feitas com calma.
\end{example}

\begin{method}[Ler a sua própria máquina]\label{met:g7:organs-at-work:read}
Uma autoexperiência com um relógio (o
\cref{met:g4:heart-and-blood:pulse} fornece a técnica de medir o
\omterm{prop:g4:heart-and-blood:pulse}{pulso}):
\begin{enumerate}
  \item registre o \omterm{prop:g4:heart-and-blood:pulse}{pulso} em repouso, sentado e calmo;
  \item faça três minutos de esforço constante --- subir e descer um
        degrau, pular corda;
  \item registre o \omterm{prop:g4:heart-and-blood:pulse}{pulso} na hora e depois a cada minuto, até ele
        voltar ao valor de repouso: é o tempo de \omterm{def:g7:organs-at-work:recovery}{recuperação};
  \item repita toda semana durante um bimestre de exercício regular:
        veja o \omterm{prop:g4:heart-and-blood:pulse}{pulso} de repouso, o \omterm{prop:g4:heart-and-blood:pulse}{pulso} de esforço e o tempo de
        \omterm{def:g7:organs-at-work:recovery}{recuperação} todos irem descendo --- a reconstrução da
        \cref{prop:g7:organs-at-work:training}, nos seus próprios
        dados.
\end{enumerate}
\end{method}

\section{Exercícios}

\begin{exercise}[$\star$]\label{exo:g7:organs-at-work:1}
Liste do que um \omterm{def:g3:bones-and-muscles:muscle}{músculo} em atividade retira mais e o que ele libera
mais.
\end{exercise}

\begin{exercise}[$\star$]\label{exo:g7:organs-at-work:2}
O que é a \omterm{prop:g7:organs-at-work:bill}{glicose}, e de onde o \omterm{def:g3:bones-and-muscles:muscle}{músculo} em atividade a tira?
\end{exercise}

\begin{exercise}[$\star$]\label{exo:g7:organs-at-work:3}
Descreva a comparação entre o \omterm{prop:g4:journey-of-food:absorb}{sangue} que entra e o que sai, e o que
ela mostra no repouso e no esforço.
\end{exercise}

\begin{exercise}[$\star$]\label{exo:g7:organs-at-work:4}
Diga as três reorganizações de abastecimento da
\cref{prop:g7:organs-at-work:supply}.
\end{exercise}

\begin{exercise}[$\star$]\label{exo:g7:organs-at-work:5}
Por que o trabalho pesado esquenta você --- e o que é tremer de frio,
nos termos deste capítulo?
\end{exercise}

\begin{exercise}[$\star$]\label{exo:g7:organs-at-work:6}
O que é a \omterm{def:g7:organs-at-work:recovery}{recuperação}, e que dívidas o corpo está acertando durante
ela?
\end{exercise}

\begin{exercise}[$\star\star$]\label{exo:g7:organs-at-work:7}
Leia o gráfico do \omterm{prop:g4:heart-and-blood:pulse}{pulso}: valor de repouso, valor do platô e mais ou
menos quanto tempo a \omterm{def:g7:organs-at-work:recovery}{recuperação} leva. O que o treino mudaria em cada
um?
\end{exercise}

\begin{exercise}[$\star\star$]\label{exo:g7:organs-at-work:8}
Explique a moleza depois do almoço e a pontada de quem corre com a
regra do redirecionamento.
\end{exercise}

\begin{exercise}[$\star\star$]\label{exo:g7:organs-at-work:9}
Por que a \omterm{def:g1:the-five-senses:sense}{pele} avermelhada do esforço é uma mudança de
\emph{entrega}, e não um defeito? Que problema ela está resolvendo, e
de quem?
\end{exercise}

\begin{exercise}[$\star\star$]\label{exo:g7:organs-at-work:10}
Liste três reconstruções pelas quais o treino levanta o teto, e o
sinal do dia a dia de cada uma.
\end{exercise}

\begin{exercise}[$\star\star$]\label{exo:g7:organs-at-work:11}
Dois alunos rodam o \cref{met:g7:organs-at-work:read}: \omterm{prop:g4:heart-and-blood:pulse}{pulsos} de
repouso $68$ e $86$; recuperações de três e de oito minutos. Que
máquina é a treinada, e o que exatamente cada um dos dois números
atesta?
\end{exercise}

\begin{exercise}[$\star\star\star$]\label{exo:g7:organs-at-work:12}
Ligue três capítulos num parágrafo só: a conta do \omterm{def:g3:bones-and-muscles:muscle}{músculo}
(\cref{prop:g7:organs-at-work:bill}), o objetivo da \omterm{def:g7:what-is-respiration:respiration}{respiração}
(\cref{prop:g7:what-is-respiration:power}) e o serviço de entregas do
\omterm{prop:g4:journey-of-food:absorb}{sangue} (\cref{prop:g4:heart-and-blood:carries}) --- como um relato
único de uma corrida só.
\end{exercise}

\section{Problema: O estudo de condicionamento da escola}

\begin{problem}[{Problema de fim de semana --- um bimestre de dados de
pulso, lidos como um fisiologista}]\label{pb:g7:organs-at-work:1}
A turma roda o \cref{met:g7:organs-at-work:read} toda segunda-feira
durante um bimestre: \omterm{prop:g4:heart-and-blood:pulse}{pulso} de repouso, \omterm{prop:g4:heart-and-blood:pulse}{pulso} depois de três minutos de
degraus e tempo de \omterm{def:g7:organs-at-work:recovery}{recuperação}. Dois registros sem identificação:
aluno A --- semana 1: repouso $84$, esforço $168$, \omterm{def:g7:organs-at-work:recovery}{recuperação} $9$
minutos; semana 12: repouso $72$, esforço $150$, \omterm{def:g7:organs-at-work:recovery}{recuperação} $5$
minutos. Aluno B --- semana 1: repouso $70$, esforço $150$,
\omterm{def:g7:organs-at-work:recovery}{recuperação} $4$ minutos; semana 12: repouso $69$, esforço $148$,
\omterm{def:g7:organs-at-work:recovery}{recuperação} $4$ minutos.

\medskip\noindent\textbf{Parte I --- Dentro de uma sessão.}
\begin{enumerate}
  \item Durante os degraus, liste o que está acontecendo com a
        retirada de \omterm{def:g7:what-is-respiration:respiration}{oxigênio} e de \omterm{prop:g7:organs-at-work:bill}{glicose} de um \omterm{def:g3:bones-and-muscles:muscle}{músculo} da coxa e com
        a produção de \omterm{def:g7:what-is-respiration:respiration}{gás carbônico} e de calor dele.
  \item Explique por que a \omterm{def:g7:what-is-respiration:respiration}{respiração} e o \omterm{prop:g4:heart-and-blood:pulse}{pulso} sobem \emph{juntos}
        durante os minutos de esforço.
  \item Um aluno se sente quente e avermelhado no terceiro minuto. Dê
        o relato da caldeira e do resfriamento.
  \item Por que o \omterm{prop:g4:heart-and-blood:pulse}{pulso} não cai para o valor de repouso no instante em
        que os degraus param?
\end{enumerate}

\medskip\noindent\textbf{Parte II --- Lendo o bimestre.}
\begin{enumerate}[resume]
  \item Descreva as mudanças de A ao longo do bimestre e diga as
        reconstruções da \cref{prop:g7:organs-at-work:training} por
        trás de cada número.
  \item Os números de B quase não se mexeram. Dê as duas leituras
        honestas --- e diga que fato a mais sobre a vida de B fora da
        aula decidiria entre elas.
  \item Na semana 5, um aluno mede logo depois de um almoço apressado
        e obtém números estranhos. Que regra da
        \cref{prop:g7:organs-at-work:supply} interferiu, e o que o
        protocolo deveria dizer sobre os horários de medição?
  \item Por que cada aluno precisa medir sempre do mesmo jeito ---
        mesmo esforço, mesmo cronômetro? (Um velho conhecido entre os
        métodos; diga o nome dele.)
\end{enumerate}

\medskip\noindent\textbf{Parte III --- A carta do
fisiologista.}
\begin{enumerate}[resume]
  \item Escreva o resumo de duas frases do bimestre de A para o mural
        da turma, usando \emph{teto} e \emph{\omterm{def:g7:organs-at-work:recovery}{recuperação}}.
  \item Uma mãe pergunta se um \omterm{prop:g4:heart-and-blood:pulse}{pulso} de repouso baixo é preocupante.
        Responda com o ciclista do
        \cref{exo:g4:heart-and-blood:11}.
  \item A turma quer um único número para acompanhar no bimestre
        seguinte. Defenda o tempo de \omterm{def:g7:organs-at-work:recovery}{recuperação} em vez do \omterm{prop:g4:heart-and-blood:pulse}{pulso} de
        esforço, usando a \cref{def:g7:organs-at-work:recovery}.
  \item Encerre com a moral do estudo numa frase: o que reconstruiu a
        máquina de A?
\end{enumerate}
\end{problem}
